\section{相关工作}
\label{sec:related_work}

\subsection{时间序列预测的深度学习方法}
时间序列研究大体经历了从统计方法（ARIMA~\cite{arima_survey}、VAR~\cite{var_survey}）到深度方法、再到 Transformer 体系的演进。围绕长序列建模效率，Informer 通过 ProbSparse 注意力降低计算复杂度~\cite{informer}；围绕结构先验注入，Autoformer 利用趋势--季节分解与自相关机制建模周期结构~\cite{autoformer}，FEDformer 将频域表示引入全局模式学习~\cite{fedformer}；iTransformer 则创新性地将传统"时间步 token 化"反转为"变量维 token 化"，以 Transformer Encoder 建模跨变量交互而非跨时间步依赖，在多变量预测任务上取得了显著提升~\cite{itransformer}。与此同时，patch 分块处理、协变量融合与线性模型等方法也持续推动性能边界。本文以 iTransformer 为骨干架构，保留其变量维 token 化与 Encoder 设计，并通过频谱条件调制增强其频率感知能力。

\subsection{频域分析与多分辨率表征}
频域方法在时间序列建模中具有重要的归纳偏置价值。FEDformer 通过傅里叶/小波变换在频域执行注意力操作~\cite{fedformer}；TimesNet 将一维序列变换为二维张量以捕获多周期模式；相关频谱分析方法利用 FFT 提取功率谱、相位谱等特征进行序列分类与异常检测~\cite{fourier_ts}。然而，多数现有方法或完全在频域内操作（引入较大架构改动），或仅将频谱作为辅助可视化工具，缺乏将其作为轻量条件信号注入通用骨干网络的机制。

多分辨率/多尺度建模通过在不同时间粒度上聚合信息以增强表征鲁棒性：金字塔结构融合多层级特征~\cite{ts_pyramid}，多尺度卷积/池化并行的思路在计算机视觉与序列任务中均有应用~\cite{multiscale_ts}。但在时序 Transformer 中，多尺度设计常伴随计算开销上升与架构复杂化等问题。本文的多分辨率频谱在概念上不同于时域多尺度池化——我们通过在不同长度的时间窗口上独立计算 FFT，获得"序列长度鲁棒"的频谱特征，尺度数量固定（3 个）且计算量可控；频带能量归一化进一步保证了不同窗口特征的可比性。

\subsection{条件调制与特征归一化}
条件调制（Conditional Modulation）通过学习输入依赖的缩放/平移参数实现样本自适应表征调整。FiLM（Feature-wise Linear Modulation）最早在视觉推理中提出~\cite{film}，后续在风格迁移、少样本学习与序列建模中得到了广泛应用~\cite{cond_norm,dyn_conv,adaptive_norm_ts}。在时间序列领域，部分工作开始探索以频率特征作为调制条件。然而，直接将 FiLM 应用于深度 Transformer 存在两个风险：(1) 调制强度过大可能淹没原始表征，破坏骨干网络的可训练性；(2) 单一调制点（仅作用于输入端）在深层网络中梯度衰减严重，调制效果逐层稀释。

本文的逐层渐进式 FiLM 直接针对上述问题：(1) 在每个 Encoder Layer 前独立注入调制，避免深层梯度衰减；(2) 调制强度按 \(\gamma_{\text{base}} \cdot \delta^{\ell}\) 指数衰减，浅层接受较强频域引导以建立全局结构感知，深层逐渐减弱以保护语义表征；(3) 可学习 gamma 与零初始化保证初始等价于无调制的骨干网络，频域信息在训练中逐步激活。该设计受稳定门控（stable gating）相关讨论的启发~\cite{stable_gating}，但将其系统化为一种有理论直觉（由粗到细）的层级化调制范式。

\subsection{非平稳性与分布漂移处理}
非平稳性是现实世界时间序列的核心挑战之一。RevIN（Reversible Instance Normalization）通过在归一化--反归一化框架下对齐每个样本的统计量，以极低的计算开销有效缓解了分布漂移对预测模型的影响~\cite{revin}，已成为许多现代时序模型的标准组件。然而 RevIN 本身不建模信号的频率特征，本文在保留 RevIN 的同时，通过 FFT 频谱条件为模型补充了归一化过程中可能被削弱的频域信息。

\subsection{科学观测序列的参数回归}
以 CEMP 恒星参数回归为代表的任务旨在从光谱/测光序列中反演 \(T_{\mathrm{eff}}\)、\(\log g\)、\([\mathrm{Fe/H}]\)、\([\mathrm{C/Fe}]\) 等物理参数。相关工作中，StarNet 采用卷积神经网络提取局部谱特征~\cite{starnet}，但在受红化、仪器差异或低信噪比影响的非平稳观测条件下，由于缺乏全局频率感知的动态调制机制，泛化鲁棒性仍有提升空间。本文统一框架通过更换任务头（预测头 \(\to\) 回归头）即可适配此类 sequence-to-vector 任务，频谱条件信号为模型提供了信噪比自适应的感知能力。

综合来看，现有工作对"多变量预测与科学观测参数回归"的统一场景，仍缺少一种同时兼顾 FFT 频谱编码、层级化自适应调制与可退化训练的轻量化框架，本文方法即在该研究空白上展开。
